\section*{Resumen}\label{abstract}

Esta memoria se centra en el estudio experimental de la viabilidad de aprovechar múltiples GPUs heterogéneas ---distintas generaciones, arquitecturas y fabricantes--- dentro de un mismo sistema, siendo lo óptimo poder utilizar todas ellas sin necesidad de descartar ninguna. El objetivo general del trabajo es determinar hasta qué punto es posible construir una máquina de cómputo de tal forma que se consiga integrar tarjetas de características distintas para así obtener un rendimiento conjunto útil. Resulta especialmente interesante emplear hardware de segunda mano e incluso \textit{legacy}, ya que es el hardware que con mayor probabilidad encontraremos en el proceso de búsqueda de dispositivos a los que darle una segunda vida útil.

Se aborda el problema mediante una metodología experimental. En primer lugar se estudian las alternativas que hacen posible la integración de múltiples tarjetas distintas, entre ellas virtualización y contenedores que permiten compartir el hardware gráfico entre múltiples entornos (VFIO y máquinas virtuales, y contenedores LXC), analizando las limitaciones técnicas de cada una. Sobre el caso de estudio formado por dos NVIDIA Tesla C1060 y una ASUS ENGTS250, de la arquitectura GT200 (2008--2009), se documentan en detalle los obstáculos encontrados: la limitación de un único driver propietario por sistema, la instalación de un driver \textit{legacy} de 2019 sobre un kernel y \textit{toolchain} de 2026 mediante parches comunitarios y DKMS, y la instalación del último toolkit CUDA compatible con estas arquitecturas (CUDA 6.5). La configuración experimental empleó dos GPUs, una Tesla C1060 y la ENGTS250, porque la tercera tarjeta no fue detectada a través del \textit{riser} PCIe. Finalmente se realiza una batería de pruebas de rendimiento (cómputo y ancho de banda), un experimento de ejecución concurrente sobre ambas GPUs y un análisis de las limitaciones del hardware en materia de telemetría energética.

Los resultados muestran que la integración de GPUs heterogéneas es viable bajo determinadas condiciones ---misma familia de driver y mismas capacidades de cómputo---, y se identifican los factores que determinan dicha viabilidad, así como las limitaciones físicas y de software que quedan abiertas como trabajo futuro.

\newpage

\section*{Abstract}

This work focuses on the experimental study of the feasibility of taking advantage of multiple heterogeneous GPUs ---different generations, architectures and vendors--- within a single system, without having to discard any of them. The general goal is to determine to what extent it is possible to build a computing machine that integrates very different graphics cards and obtains useful aggregate performance, mainly using second-hand or legacy hardware.

The problem is approached through an experimental methodology. First, the virtualization and container alternatives that allow sharing graphics hardware between multiple environments (VFIO with virtual machines, and LXC containers) are studied, analysing the technical limitations of each one. On the case study formed by two NVIDIA Tesla C1060 and an ASUS ENGTS250, both from the GT200 architecture (2008--2009), the obstacles found are documented in detail: the single proprietary driver limitation per system, the installation of a 2019 legacy driver on a 2026 kernel and toolchain by means of community patches and DKMS, and the installation of the last CUDA toolkit compatible with these architectures (CUDA 6.5). The experimental configuration used two GPUs, one Tesla C1060 and the ENGTS250, because the third card was not detected through the PCIe riser. Finally, a set of performance benchmarks (compute and memory bandwidth), a concurrent execution experiment on both GPUs, and an analysis of the hardware limitations regarding energy telemetry are carried out.

The results show that the integration of heterogeneous GPUs is feasible under certain conditions ---same driver family and same compute capabilities--- and the factors that determine such feasibility are identified, as well as the physical and software limitations that remain open as future work.
